\documentclass{article}
\usepackage[utf8]{inputenc}
\usepackage{graphicx}
\usepackage[left=2cm,right=2cm,top=2cm,bottom=2cm]{geometry}
\usepackage{amsmath,amssymb,amsfonts,amsthm,dsfont, braket}
\usepackage{blindtext}
\usepackage{multicol}
\title{Moltbook Network Analysis}
\author{Manuel Antonio Pianini, Davide Usai, Alice Vistalli}
\date{\today}
\begin{document}
\maketitle
\newpage
\tableofcontents
\newpage
\setcounter{page}{1}
\begin{multicols}{2}
   \section*{Scelta del dataset}
   Il network oggetto del progetto è la rete di connessioni tra agenti AI del social Moltbook, di cui indaghiamo,
   oltre all'analisi convenzionale delle proprietà del grafo statico (Chap.\ref{chap_analisi}) e dinamico (Chap. \ref{chap_analisi dinamic}),
   l'organizzazione in community (Chap.\ref{chap_community detection}).
   Ci interroghiamo anche su alcuni aspetti dell'evoluzione delle community, evidenziando le peculiarità della specifica rete.
   I dati utilizzati sono pubblici ed accessibili da huggingface/SimulaMet/moltbook-observatory-archive (previa registrazione ad HuggingFace),
   ma onerosi in memoria: selezioniamo solo alcuni dataset consultabili, riducendoci ad agenti, post e commenti, da cui estrarre il grafo utile lavorando con una mole
   gestibile di dati.
\end{multicols}

\hrulefill

\begin{multicols}{2}
   \section{Data Collection}
   I tre dataset scelti sono dizionari della forma
   \begin{itemize}
      \item agenti: 'id', 'name', 'description', 'karma', 'follower\_count', 'following\_count', 'is\_claimed',
         'owner\_x\_handle', 'first\_seen\_at', 'last\_seen\_at', 'created\_at', 'avatar\_url', 'dump\_date'
      \item post: 'id', 'agent\_id', 'agent\_name', 'submolt', 'title', 'content', 'url', 'score', 'comment\_count', 
         'created\_at', 'fetched\_at', 'is\_pinned', 'dump\_date'
      \item commenti: 'id', 'post\_id', 'agent\_id', 'agent\_name', 'parent\_id', 'content', 'score', 'created\_at',
         'fetched\_at', 'dump\_date'
   \end{itemize}
   Nella nostra analisi, ci avvarremo principalmente degli ID di agenti, post e commenti, del nome degli agenti e dei timestamps.
   La scelta di ignorare il resto sarà motivata nel proseguo dell'analisi.

   \subsection{Coerenza e qualità di dati}
   I dataset vengono aggiornati separatamente dal provider, pertanto non sono perfettamente coerenti: può
   succedere ad esempio che un agente si disiscriva dal social, venendo depennato dal dataset \textit{agenti}
   ma ciò non cancella i suoi eventuali tracce nei dataset \textit{post} e \textit{commenti}.
   Questo porta alla comparsa alcuni \textit{edges fantasma} di difficile gestione.
   
   Al fine di creare un grafo consistente, definiamo come edge l'\textbf{interazione tra agenti}, ossia l'atto
   di commentare un post o il commento di un altro agente o se stesso, mentre un nodo rappresenterà un
   \textbf{agente} con un accezione ristretta: nel solo caso in cui questo abbia postato ed il suo post sia stato commentato.
   Questa scelta induce alcune notevoli conseguenze:
   \t ogni nodo ha almeno un edge, perciò compariranno i nodi singoli (disconnessi) solo se questi hanno
   self-loops ( agenti che si sono commentati da soli);
   \t possiamo ricavare la lista degli edges e dedurne quella dei nodi, unendo i dizionari con la funzione
   \textit{merging} di pandas ed usando come dataset principale quello dei commenti, accoppiando gli ID degli agenti che hanno interagito;
   \t questi edges sono per natura direzionati e pesati (nel caso di molteplici interazioni).
   
   \subsection{Direzione e peso degli edges}
   Dopo esserci ristretti ai soli autori commentati, a questo punto, la lista degli edges è un insieme di
   coppie di ID degli agenti, con un \textit{former author} ed un \textit{latter author}.
   Questa ridefinizione è il modo più comodo per gestire il caso di un commento ai commenti: se ciò avviene in un caso, infatti potrebbe avvenire all'infinito.
   Fatto ciò, il numero di coppie della lista è pari al numero dei commenti.
   All'interno di una stessa coppia, ordiniamo alfanumericamente gli ID di \textit{authors}: in questo modo
   \[ \textit{agente B}\rightarrow\textit{agente A} \Rightarrow \textit{agente A}\rightarrow\textit{agente B} \]
   \[ \textit{agente A}\rightarrow\textit{agente B} \Rightarrow \textit{agente A}\rightarrow\textit{agente B} \]
   il grafo perde direzionalità.
   Riordinando poi l'intera lista possiamo contare (eventualmente segnando i pesi) gli edges effettivi e non direzionati, raggruppando le coppie identiche.
\end{multicols}

\hrulefill

\section{Network Analysis}
Il grafo costruito a questo punto viene analizzato con il pacchetto \textit{networkx} di Python.

Il grafo Analysis of the graph data using NetworkX and Matplotlib. The code reads a CSV file containing graph data,
constructs a graph from it, and visualizes the graph using a spring layout.

\end{document}

MAKE A WEIGHTED VERSION OF DATASET, UN-TIMESTAMPED. CHECK IF WEIGHTS SIGNIFICANTLY CHANGE THE RESULTS

STATIC GRAPH ANALYSIS:
Degree distribution, V
avg degree (---> small world? scale free?),        
Connected components analysis;                              
Path analysis (avg path lenght, diameter, )
Clustering Coefficient (global + local, avg),     V         ----> low clustering? ---> barabasi-albert Clust_BA = (ln N^2)/N = 0.0052 for N=3111
Density analysis (sparse? dense? self-interactions?);       ---->unweighted density = 0.0024, weighted density = 0.0009??
Centrality analysis (degree,
    connectivity: katz?,
    geometrical: closeness? farness? betweeness?)
S-analysis with edge weight?

CLUSTER 1: COMMUNITY DISCOVERY
Assortativity (homophily, degree correlation)                   LOUVAIN ha senso qui?
Bridges detection (GIRVAN-NEUMANN) and neighbourhood
k-cliques?
Label propagation?


and many others ...
Compare with random models ER and BA, and see how the properties differ.
RANDOM MODEL: avg distance and degree in ER is it sparse?
BARABASI: max degree of a hub is consistent with osservation?


DYNAMIC ANALYSIS: is it present some sort of preferenctial attachment? or the new nodes and links are random?
'''
